\documentclass[preview]{standalone}
\usepackage[english]{babel}
\usepackage{amsmath}
\usepackage{amssymb}
\begin{document}
\begin{center}
\special{dvisvgm:raw <g id='unique000'>}Su ogni cella la funzione è vicina alla sua approssimazione lineare, a meno di un errore $< \varepsilon$. \\Questo perché $f$ è olomorfa su $\Omega$, per ipotesi, e quindi \\il limite del rapporto incrementale converge alla derivata prima \\uniformemente in ogni cella quadrata\special{dvisvgm:raw </g>}
\end{center}
\end{document}